\chapter{Mô Hình Quy Trình Phần Mềm và Kế Hoạch Dự Án}
\label{ch:quytrinh_duan}
\textit{Vo Trong Nguyen}

\section{Lựa Chọn Mô Hình Quy Trình Phần Mềm (Process Model)}
Dựa trên các nguyên lý hướng dẫn tại Chương 3 và Chương 4 của giáo trình Pressman, hệ thống định hướng nghề nghiệp này áp dụng mô hình \textbf{Agile/Incremental (Scrum-based)}. Nền tảng đòi hỏi sự tích hợp liên tục của các thuật toán AI và phân tích xu hướng công nghệ ngoài thị trường lao động vốn thay đổi hằng ngày, do đó việc phát triển theo các phân đoạn tăng trưởng (Increments) giúp giảm thiểu rủi ro lỗi thời của hệ thống.

\section{Tổ Chức Sprint và Kế Hoạch Phân Phối Nhân Sự}
Dự án được phân rã thành các chu kỳ Sprint cố định kéo dài 2 tuần. Phân chia vai trò cụ thể trong đội ngũ phát triển bao gồm:
\begin{itemize}
	\item \textbf{Thành viên 1:} Scrum Master, quản lý tiến độ tổng quan và rủi ro.
	\item \textbf{Thành viên 2 \& 3:} Kỹ sư phân tích yêu cầu (Business Analyst) chịu trách nhiệm xây dựng tài liệu đặc tả SRS và mô hình Use Case hệ thống.
	\item \textbf{Thành viên 4:} Kiến trúc sư hệ thống (System Architect), phụ trách cấu trúc mã nguồn nền tảng FastAPI Backend.
	\item \textbf{Thành viên 5:} Nhà phát triển Front-End (UI Developer), phụ trách hiện thực hóa giao diện Next.js dạng các cấu trúc component có khả năng tái sử dụng.
	\item \textbf{Thành viên 6:} Kỹ sư DevOps \& Cấu hình, thiết lập Git Workflow, Jira Board liên kết mã token môi trường và chuyển đổi ERD sang định dạng Markdown.
	\item \textbf{Thành viên 7:} Kỹ sư dữ liệu (Database Designer), cài đặt SQL Server vật lý và mô hình hóa Sơ đồ hoạt động hệ thống.
\end{itemize}

\section{Phân Tích và Giảm Thiểu Rủi Ro Dự Án}
Theo khung lý thuyết quản lý rủi ro của Pressman, dự án đối mặt với rủi ro kỹ thuật lớn nhất nằm ở độ trễ xử lý của thuật toán kết xuất lộ trình học tập bằng trí tuệ nhân tạo. Giải pháp giảm thiểu là xây dựng một bộ Core-Engine tối giản (MVP) ngay trong Sprint đầu tiên để kiểm thử hiệu năng.